\documentclass{amsart}
% For dealing with references we use the comment environment
\usepackage{verbatim}
\newenvironment{reference}{\comment}{\endcomment}
%\newenvironment{reference}{}{}
\newenvironment{slogan}{\comment}{\endcomment}
\newenvironment{history}{\comment}{\endcomment}

% For commutative diagrams we use Xy-pic
\usepackage[all]{xy}

% We use 2cell for 2-commutative diagrams.
\xyoption{2cell}
\UseAllTwocells

% We use multicol for the list of chapters between chapters
\usepackage{multicol}

% This is generally recommended for better output
\usepackage{lmodern}
\usepackage[T1]{fontenc}

% For cross-file-references

% Package for hypertext links:
\usepackage{hyperref}

% For any local file, say "hello.tex" you want to link to please

% Theorem environments.
%
\theoremstyle{plain}
\newtheorem{theorem}[subsection]{Theorem}
\newtheorem{proposition}[subsection]{Proposition}
\newtheorem{lemma}[subsection]{Lemma}

\theoremstyle{definition}
\newtheorem{definition}[subsection]{Definition}
\newtheorem{example}[subsection]{Example}
\newtheorem{exercise}[subsection]{Exercise}
\newtheorem{situation}[subsection]{Situation}

\theoremstyle{remark}
\newtheorem{remark}[subsection]{Remark}
\newtheorem{remarks}[subsection]{Remarks}

\numberwithin{equation}{subsection}

% Macros
%
\def\lim{\mathop{\mathrm{lim}}\nolimits}
\def\colim{\mathop{\mathrm{colim}}\nolimits}
\def\Spec{\mathop{\mathrm{Spec}}}
\def\Hom{\mathop{\mathrm{Hom}}\nolimits}
\def\Ext{\mathop{\mathrm{Ext}}\nolimits}
\def\SheafHom{\mathop{\mathcal{H}\!\mathit{om}}\nolimits}
\def\SheafExt{\mathop{\mathcal{E}\!\mathit{xt}}\nolimits}
\def\Sch{\mathit{Sch}}
\def\Mor{\mathop{\mathrm{Mor}}\nolimits}
\def\Ob{\mathop{\mathrm{Ob}}\nolimits}
\def\Sh{\mathop{\mathit{Sh}}\nolimits}
\def\NL{\mathop{N\!L}\nolimits}
\def\CH{\mathop{\mathrm{CH}}\nolimits}
\def\proetale{{pro\text{-}\acute{e}tale}}
\def\etale{{\acute{e}tale}}
\def\QCoh{\mathit{QCoh}}
\def\Ker{\mathop{\mathrm{Ker}}}
\def\Im{\mathop{\mathrm{Im}}}
\def\Coker{\mathop{\mathrm{Coker}}}
\def\Coim{\mathop{\mathrm{Coim}}}

% Boxtimes
%
\DeclareMathSymbol{\boxtimes}{\mathbin}{AMSa}{"02}

%
% Macros for moduli stacks/spaces
%
\def\QCohstack{\mathcal{QC}\!\mathit{oh}}
\def\Cohstack{\mathcal{C}\!\mathit{oh}}
\def\Spacesstack{\mathcal{S}\!\mathit{paces}}
\def\Quotfunctor{\mathrm{Quot}}
\def\Hilbfunctor{\mathrm{Hilb}}
\def\Curvesstack{\mathcal{C}\!\mathit{urves}}
\def\Polarizedstack{\mathcal{P}\!\mathit{olarized}}
\def\Complexesstack{\mathcal{C}\!\mathit{omplexes}}
% \Pic is the operator that assigns to X its picard group, usage \Pic(X)
% \Picardstack_{X/B} denotes the Picard stack of X over B
% \Picardfunctor_{X/B} denotes the Picard functor of X over B
\def\Pic{\mathop{\mathrm{Pic}}\nolimits}
\def\Picardstack{\mathcal{P}\!\mathit{ic}}
\def\Picardfunctor{\mathrm{Pic}}
\def\Deformationcategory{\mathcal{D}\!\mathit{ef}}

\setlength{\emergencystretch}{3em}
\begin{document}
\title[Stacks Project, Tag 0CF0]{635 --- Invariant open subsets under elementary affine automorphisms}
\maketitle
\noindent Copyright (C) 2005--2025 Johan de Jong. Stacks Project authors. Source: \href{https://stacks.math.columbia.edu/tag/0CF0}{Tag 0CF0}.

\noindent Editorial difficulty note (not author text). Difficulty H2: The proof distinguishes characteristic zero from positive characteristic and proves Zariski-density of transformation orbits using transcendental coordinates and polynomial translations. The resulting exceptional finite-point classification requires a substantive algebraic-geometric orbit argument..

\noindent Editorial provenance note (not author text). History: excerpt from revision \texttt{a0444\allowbreak 6e57e\allowbreak c1fbc\allowbreak 252a8\allowbreak 71afc\allowbreak ec775\allowbreak 2fb28\allowbreak 07b14}, descent.tex, lines 7629--7684. Reference presentation was adapted; original-author context and core support blocks were retained or added. The mathematical wording is unchanged. Licensed under GNU FDL 1.2 or later; no invariant sections or cover texts. See ../licenses/COPYING. Context blocks reproduce author definitions and supporting results; they are not additional counted problems.

\begin{lemma}
\label{lemma-orbits}
Let $k$ be a field. Let $n \geq 2$. For $1 \leq i, j \leq n$ with
$i \not = j$ and $d \geq 0$ denote $T_{i, j, d}$ the automorphism
of $\mathbf{A}^n_k$ given in coordinates by
$$
(x_1, \ldots, x_n) \longmapsto
(x_1, \ldots, x_{i - 1}, x_i + x_j^d, x_{i + 1}, \ldots, x_n)
$$
Let $W \subset \mathbf{A}^n_k$ be a nonempty open subscheme
such that $T_{i, j, d}(W) = W$ for all $i, j, d$ as above.
Then either $W = \mathbf{A}^n_k$ or the characteristic of $k$
is $p > 0$ and $\mathbf{A}^n_k \setminus W$ is a finite set
of closed points whose coordinates are algebraic over $\mathbf{F}_p$.
\end{lemma}

\begin{proof}
We may replace $k$ by any extension field in order to prove this.
Let $Z$ be an irreducible component of $\mathbf{A}^n_k \setminus W$.
Assume $\dim(Z) \geq 1$, to get a contradiction.
Then there exists an extension field $k'/k$ and a $k'$-valued
point $\xi = (\xi_1, \ldots, \xi_n) \in (k')^n$ of
$Z_{k'} \subset \mathbf{A}^n_{k'}$
such that at least one of $x_1, \ldots, x_n$ is transcendental over the
prime field. Claim: the orbit of $\xi$ under the group generated by
the transformations $T_{i, j, d}$ is Zariski
dense in $\mathbf{A}^n_{k'}$. The claim will give the desired contradiction.

\medskip\noindent
If the characteristic of $k'$ is zero, then already the operators
$T_{i, j, 0}$ will be enough since these transform $\xi$ into
the points
$$
(\xi_1 + a_1, \ldots, \xi_n + a_n)
$$
for arbitrary $(a_1, \ldots, a_n) \in \mathbf{Z}_{\geq 0}^n$.
If the characteristic is $p > 0$, we may assume after renumbering
that $\xi_n$ is transcendental over $\mathbf{F}_p$. By
successively applying the operators $T_{i, n, d}$ for
$i < n$ we see the orbit of $\xi$ contains the elements
$$
(\xi_1 + P_1(\xi_n), \ldots, \xi_{n - 1} + P_{n - 1}(\xi_n), \xi_n)
$$
for arbitrary $(P_1, \ldots, P_{n - 1}) \in \mathbf{F}_p[t]$.
Thus the Zariski closure of the orbit contains the coordinate
hyperplane $x_n = \xi_n$. Repeating the argument with a different
coordinate, we conclude that the Zariski closure contains
$x_i = \xi_i + P(\xi_n)$ for any $P \in \mathbf{F}_p[t]$
such that $\xi_i + P(\xi_n)$ is transcendental over $\mathbf{F}_p$.
Since there are infinitely many such $P$ the claim follows.

\medskip\noindent
Of course the argument in the preceding paragraph also applies
if $Z = \{z\}$ has dimension $0$ and the coordinates of $z$
in $\kappa(z)$ are not algebraic over $\mathbf{F}_p$. The lemma follows.
\end{proof}
\end{document}
